\uuid{cooY}
\titre{Matrice/Somme/Produit}
\auteur{hanani ; mzari}
\organisation{exo7}
\licence{CC BY-NC-SA 4.0}
\datecreate{2020-01-20}
\niveau{L1}
\module{Algèbre}
\chapitre{Matrice}
\sousChapitre{Propriétés élémentaires, généralités}
\difficulte{1}
\theme{Calcul matriciel, Matrice, Somme, Produit}
\exo7qcmid{lille-2.matrices.001}

\contenu{
\begin{question}
\qclassification{108.01}
\qtags{motcle=Matrice/Somme/Produit}
Soit $A$ et $B$ deux matrices. Quelles sont les assertions vraies ?
\begin{answers}
\good{Si la matrice $A+B$ est définie, alors $B+A$ est définie}.
\bad{Si la matrice $A+B$ est définie, alors $AB$ est définie}.
\bad{Si la matrice $AB$ est définie, alors $BA$ est définie}.
\good{Si la matrice $A+B$ est définie, alors $A^tB$ est définie, où $^tB$ est la transposée de la matrice $B$}.
\end{answers}
\begin{explanations} La somme de deux matrices est définie si les deux matrices admettent la même taille. Dans ce cas, on a $A+B=B+A$. Le produit $AB$ de deux matrices est défini si le nombre de colonnes de $A$ est égal au nombre de lignes de $B$. Le produit n'est pas commutatif.
\end{explanations}
\end{question}
}
